
\section{Proof sketch}
\begin{frame}
\textbf{Goal: }Upper bound $\text{spar}(g)$ in terms of $\text{rank}^{1/3} (g \circ \text{AND}_2)$

\begin{itemize}
	\item<1-> Assume, for the sake of contradiction, $\text{spar}(g)$ is large and $\text{rank}^{1/3} (g \circ \text{AND}_2)$ is small.
	\item<2-> From the second assumption, we have functions
	$$ a_i, b_i : \{0,1\}^n \rightarrow [0,1] , i= 1, 2, \cdots , m$$
	and constants $\alpha_i, i = 1, 2, \cdots , m$ such that
	
	$$\left| g(x \land y) - \displaystyle \sum_{i=1}^{m} \alpha_i a_i(x) \land b_i(y)\right| \leq 1/3$$
	\item<2-> By standard techniques (eg John's lemma), we can assume $\sum |\alpha_i| \leq L$ (where $\log L \leq \text{polylog} (m)$)
\end{itemize}

\end{frame}

\begin{frame}
	\textbf{Goal: }Upper bound $\text{spar}(g)$ in terms of $\text{rank}^{1/3} (g \circ \text{AND}_2)$
	
	\begin{itemize}
		\item<1-> We shall construct a function $f$ (to be obtained from $g$) such that
		\begin{itemize}
			\item Approximate degree of $f$ is small
			\item Degree of $f$ is high
		\end{itemize}
		This contradicts known relations between degree and approximate degree (first proven by Nisan-Szegedy; constants were optimized later).
		
		\item<2-> We shall get a low-degree approximating polynomial from the inequality 
			$$\left| g(x \land y) - \displaystyle \sum_{i=1}^{m} \alpha_i a_i(x) \land b_i(y)\right| \leq 1/3$$
			whereas the high sparsity of $g$ will guarantee the original polynomial has high degree.
	\end{itemize}
	
\end{frame}


\begin{frame}
	\textbf{Goal: }Upper bound $\text{spar}(g)$ in terms of $\text{rank}^{1/3} (g \circ \text{AND}_2)$
	
	\begin{itemize}
		\item<1-> Note that every function $f: \{0,1\}^n \rightarrow \mathbb{R}$ can also be represented as a function $\tilde{f}: \{ \pm 1\}^n \rightarrow \mathbb{R}$. It is easy to see that their degrees are the same. \newline
		
		 \textbf{Convention: }For a function $f: \{0,1\}^n \rightarrow \mathbb{R}$, we denote its equivalent function with the domain $\{ \pm 1\}^n$ by $\tilde{f}: \{ \pm 1\}^n \rightarrow \mathbb{R}$.
		\item<2-> For upper-bounding the approximate degree, we shall look at functions in the $\{\pm 1\}^n$ domain. For lower bounding the degree, we shall look at functions in the $\{ 0, 1\}^n$ domain. Since the degree is invariant under change of domain, this is ok.
	\end{itemize}
	
\end{frame}

\begin{frame}
	
	\begin{itemize}
		\item \textbf{Observation: }Given the expression 
		$$ g(x \land y) \approx \displaystyle \sum_{i=1}^{m} \alpha_i a_i(x) \land b_i(y),$$
		we can hit each $a_i$ by a restriction+averaging operator $(\rho, \delta)$, replace each $b_i$ by a constant, and get an approximation for $g_{\rho}$
	\end{itemize}
\end{frame}


\begin{frame}
	
	\textbf{Restriction+averaging operators: }
	
	A \textbf{restriction+averaging operator} is denoted by a string $R \in \{0,1,*,\Delta\}^n$. It converts a function $f: \{0,1\}^n \rightarrow \mathbb{R}$ to a function $f: \{0,1\}^{\text{free}(R)} \rightarrow \mathbb{R}$ where $\text{free}(R)= \{i | R_i=*\}$. The transformation is as follows:
	\begin{itemize}
		\item If $R_i=b \in \{0,1\}$, fix the $i$-th coordinate to $b$.
		\item If $R_i= \Delta$, average over the $i$-th coordinate.
		\item If $R_i=*$, keep the $i$-th coordinate free.
	\end{itemize}
	
	We shall associate a partial assignment $\rho \in \{0,1,*\}^n$ to each restriction+averaging operator $R$: $\rho$ is obtained from $R$ by changing all the $\Delta$'s to 0's. 
	
	
\end{frame}

\begin{frame}
	\textbf{Goal: }Upper bound $\text{spar}(g)$ in terms of $\text{rank}^{1/3} (g \circ \text{AND}_2)$
	
	\begin{itemize}
		\item<1-> We shall define a distribution $\mathcal{D}$ over \textit{restriction+averaging operators} $R$ such that for any $\tilde{a}: \{ \pm 1\}^n \rightarrow [0,1]$,
		$$ \text{E}_{R \sim \mathcal{D}} [ |\tilde{a}_{R}|^{\geq d}] \leq \varepsilon$$
		
		where $|\tilde{f}|^{\geq d}$ denotes $\ell^1$ Fourier weight above level $d$
		\item<2-> ... and such that $\text{E}_{R \sim \mathcal{D}}[ f_{\rho}] \geq D$, where $\rho$ is the partial assignment associated with $R$.
	\end{itemize}	
	
\end{frame}

\begin{frame}

	\begin{itemize}
		\item<1-> The first property, along with the expression for approximate rank, will give that $\text{approx-deg}(g_{\rho}) \leq d$ for a typical $R \sim \mathcal{D}$
		\item<2-> The second property will give that $\text{deg}(g_{\rho}) \geq D$ for a typical $R \sim \mathcal{D}$.
		\item<3-> If the gap between $d$ and $D$ is large enough, we are done.
	\end{itemize}	
	
\end{frame}


\begin{frame}{Obtaining an approximation for $f_{\rho}$}
\begin{itemize}
	\item<1-> Recall the expression we got from the low-approximate rank assumption:
	$$ g(x \land y) \approx \displaystyle \sum_{i=1}^{m} \alpha_i a_i(x) b_i(y)$$
	$$ a_i, b_i: \{0,1\}^n \rightarrow [0,1]$$
	$$ \displaystyle \sum_i |\alpha_i| \leq L$$
	\item<2-> \textbf{Claim: }Given a restriction+averaging operator $R$, with $\rho$ its associated partial assignment, there exist constants $c_1, c_2, \cdots , c_m \in [0,1]$ such that
	$$ g_{\rho} (z) \approx \displaystyle \sum_{i=1}^{m} \alpha_i c_i (a_i)_{R} (z)$$

\end{itemize}	

\end{frame}

\begin{frame}{Obtaining an approximation for $f_{\rho}$}
	\begin{itemize}
		\item<1-> \textbf{Claim: }Given a restriction+averaging operator $R$, with $\rho$ its associated partial assignment, there exist constants $c_1, c_2, \cdots , c_m \in [0,1]$ such that
		$$ g_{\rho} (z) \approx \displaystyle \sum_{i=1}^{m} \alpha_i c_i (a_i)_{R} (z)$$
		\item<2-> I will show this only for the special case $R= \{\Delta , *, *, *, \cdots , *\}$. General case should be clear from this.
		\item<3-> So, $\rho = \{0, *, *, *, \cdots , *\}$.
	\end{itemize}	
	
\end{frame}


\begin{frame}{Obtaining an approximation for $f_{\rho}$}
	\begin{itemize}
		\item<1->I will show this only for the special case $R= \{\Delta , *, *, *, \cdots , *\}$. General case should be clear from this. So, $\rho = \{0, *, *, *, \cdots , *\}$.
		\item<2-> Plugging $y_0= (0, 1, 1, \cdots , 1), x= (1, z_2, z_3, \cdots , z_n)$, we get
		$$ g(0, z_2, z_3, \cdots , z_n) \approx \displaystyle \sum_{i=1}^{m} \alpha_i a_i(1, z_2, \cdots , z_n) b_i(y_0)$$
		Plugging $y_0= (0, 1, 1, \cdots , 1), x= (0, z_2, z_3, \cdots , z_n)$,
		$$ g(0, z_2, c_3, \cdots , z_n) \approx \displaystyle \sum_{i=1}^{m} \alpha_i a_i(0, z_2, z_3, \cdots , z_n) b_i(y_0) $$
		\item<3-> Average these two expressions:
		$$ g(0, z_2, \cdots , z_n) \approx \displaystyle \sum_{i=1}^{m} \alpha_i b_i(y_0) (a_i)_R(z)$$
			
	\end{itemize}	
	
\end{frame}

\begin{frame}{Exhibiting a gap between degree and approximate degree}
	\begin{itemize}
		\item<1-> \textbf{Conclusion: }If $R$ is a restriction+averaging operator and $\rho$ is the associated partial assignment, we get an approximation for $g_{\rho}$ in terms of the $(a_i)_R$'s.
		\item<2-> \textbf{Our next goal: }Using the fact that sparsity of $g$ is high, we define a distribution $\mathcal{D}$ over restriction+averaging operators $R$ such that:
		\begin{itemize}
			\item For all $a: \{0,1\}^n \rightarrow [0,1]$ ,$\text{E}_{R \sim \mathcal{D}}[\tilde{\text{deg}}(a_{R})] \leq d$ \textcolor{red}{(we will do something slightly different actually)}
			\item $\text{E}_{(R \sim \mathcal{D})} [\text{deg}(g_{\rho})] \geq D$ where $\rho$ is the partial assignment corresponding to $R$.
		\end{itemize} 
		\item<3-> If we can do this with a sufficiently big gap between $d$ and $D$, from the previous observation we get that the approximate degree of $g_{\rho}$ is very small; whereas teh second property implies the exact degree of $g_{\rho}$ is very large. But Nisan-Szegedy implies for any Boolean function, degree and approximate degree are polynomially related. Contradiction.
		
		
	\end{itemize}	
	
\end{frame}



\begin{frame}{Simplifying any bounded function}
\textbf{First requirement: }
\begin{itemize}
\item<1-> This is what we will show: \newline
For any $\tilde{a}: \{ \pm 1\}^n \rightarrow [0,1]$, 
$$ \text{E}_{R \sim \mathcal{D}} [|| \tilde{a}_R ||^{\geq d}] \leq \varepsilon$$
where $||f||^{\geq d}$ denotes total $\ell^1$ Fourier mass above level $d$. \newline
Notice that we switched to $\{ \pm 1\}$ domain here.
\item<2-> If we achieve this, the total $\ell^1$ Fourier mass above level $d$ of the function
$$ \displaystyle \sum_{i=1}^{L} \alpha_i c_i (a_i)_R$$ will be at most $L \varepsilon$, which we will set to be $<1/100$. Thus, on expectation, the approximate degree of $g_{\rho}$ will be $\leq d$.
\end{itemize}
\end{frame}

\begin{frame}{Simplifying any bounded function}
\textbf{Why switching to the $\{ \pm 1 \}$ basis is helpful:}
\begin{itemize}
	\item<1-> Let $\tilde{f}: \{ \pm 1\}^n \rightarrow \mathbb{R}$ be any function. Its Fourier expansion is
	$$ \tilde{f} = \displaystyle \sum_{S \subseteq [n]} \hat{\tilde{f}}(S) \chi_S$$
	\item<2-> Suppose we hit $\tilde{f}$ by a restriction+averaging operator $R$, which averages over some subset $T$ of variables. Then, we have
	$$ \tilde{f}_R = \displaystyle \sum_{S \cap T = \phi} \hat{\tilde{f}}(S) \chi_S$$
	\item<3->Thus, averaging out over a variable kills all monomials containing it.
\end{itemize}
\end{frame}

\begin{frame}{Defining the distribution}

\begin{itemize}
	\item<1-> Assume $g: \{0,1\}^n \rightarrow \{0,1\}$ has $\geq s$ non-zero monomials in the deMorgan basis. Let $\mathcal{F} \subseteq 2^{[n]}$ denote the set of nonzero monomials. 
	\item<2-> By \textbf{Sauer-Shelah lemma}, the VC dimension of $\mathcal{F}$ is at least $\Omega (\log(s)/\log(n))$.
	\item<3-> Let $T \subseteq [n]$ be a shattered set of size $\geq \Omega(\log(s)/\log(n))$.
	\item<4-> For any $A \subseteq T$, there exists $S \in \mathcal{F}$ such that $S \cap T = A$.
\end{itemize}

\end{frame}


\begin{frame}{Defining the distribution}
	
	\begin{itemize}
		\item<1-> Assume $g: \{0,1\}^n \rightarrow \{0,1\}$ has $\geq s$ non-zero monomials in the deMorgan basis. Let $\mathcal{F} \subseteq 2^{[n]}$ denote the set of nonzero monomials. 
		\item<2-> By \textbf{Sauer-Shelah lemma}, the VC dimension of $\mathcal{F}$ is at least $\Omega (\log(s)/\log(n))$.
		\item<3-> Let $T \subseteq [n]$ be a shattered set of size $\geq \Omega(\log(s)/\log(n))$.
		\item<4-> For any $A \subseteq T$, there exists $S \in \mathcal{F}$ such that $S \cap T = A$.
	\end{itemize}
	
\end{frame}

\begin{frame}{Defining the distribution}
	
	\begin{itemize}
		\item<1-> Let $T$ be the large shattered subset.
		\item<2-> For every $i \in T$, set 
		$$
		R(i) = \begin{cases}
			* & \text{ with probability }p \\
			\Delta & \text{ with probability }1-p
		\end{cases}
		$$
	(Here $p$ is small but non-negligible.)
	\item<3-> Let $A = \{i | i \in T, R(i)=*\}. $Look at the polynomial
	$$ g_{A,T} = \displaystyle \sum_{S : S \cap T = A } \alpha_S x_{S \setminus A}$$
	This is defined on $\{0,1\}^{[n] \setminus T}$, and is nonzero since $T$ is shattered. So there exists a $y \in \{0,1\}^{[n] \setminus T}$ such that $g_{A,T}(y) \neq 0$. 
	\item<4-> Set $R(i)=y_i$ for $i \not \in T$.
	\end{itemize}
\end{frame}

